\section{Limitations}

\paragraph{Controlled stress and artifact-bounded contracts.}
The evidence comes from controlled custom stress sets, a local mock pre-commit contract, an independently specified held-out contract, and sandboxed/replayed traces. It does not establish safety for real SaaS, browser, banking, email, calendar, file-sharing, Slack, or workspace systems, and the numbers should not be interpreted as deployed false-allow or false-denial rates. The experiments avoid real external side effects and do not model post-commit rollback.

\paragraph{Synthetic and deterministic artifacts.}
The E55/E57 local contract is deterministic and synthetic. Its resource schemas, aliases, and authorization context are designed for the experiment. E60 adds an independently specified contract and E61 adds sandboxed realistic trace replay plus a saved external replay subset, but these remain generated/replayed artifact evaluations rather than live third-party deployments. The external subset uses trace metadata and rule-derived sidecar atoms rather than independent human annotation. E57-v2 perturbation and independent reference-authorizer checks reduce lexical-template and implementation-coupling concerns; they do not prove general transfer.

\paragraph{Authorization-system scope.}
The E55-v2 prototype assumes that the local environment can provide typed authorization context, canonical resource information, and provenance/control-source summaries. It is therefore not a complete enterprise permission system and does not handle malicious tool implementations, compromised operating systems, browser/SaaS compromise, or policy lifecycle management. The contribution is the pre-commit semantic interface, not a replacement for production authorization infrastructure.

\paragraph{Baseline scope.}
Released checkpoints are evaluated on our custom stress interface. These results are useful for diagnosing binding behavior under our stress axes, but they are not original-paper benchmark reproductions. We therefore avoid claims that one external system is globally safer or weaker than another.

\paragraph{Human audit findings.}
The human audit found corrections. This is a strength for claim hygiene, but a limitation for label certainty. The paper should use corrected-label sensitivity wording and should not claim that the audit perfectly confirmed all labels, atoms, or violation reasons.

\paragraph{Remaining evaluation gaps.}
E60--E64 address the largest prior gaps by adding held-out contracts, realistic trace replay, extraction decomposition, interface-burden analysis, and stronger comparable baselines. Remaining gaps are live integration with independently maintained authorization systems, externally authored and human-labeled trace corpora, uncertainty estimates over larger domains, and released-system adapters that run their original protocols end to end. These would test whether the atom interface remains useful outside the controlled artifact.
